%%% lAtEx tEMPLATE: Article/Thesis/etc. with colored headings and special fonts
%%%
%%% Source: http://www.howtotex.com/
% vim: set spell spelllang=es syntax=tex :

\documentclass[12pt]{article}
\usepackage[T1]{fontenc}
\usepackage{styles/apuntes-estilo}
\usepackage{styles/egyptian}
\usepackage{fancyhdr,lastpage}
\usepackage{hyperref}
\usepackage[inline]{enumitem}
\usepackage{pmboxdraw} % Para poder usar la salida del comando tree

\def\maketitle{

\makeatletter{
    \color{blue} \centering \huge \sc
    \textbf{
        Segundo examen parcial\\
        \small \vspace*{-8pt} \color{black}
        Recuperatorio
        \vspace*{8pt}
    }\\
    \par
}
\makeatother

\makeatletter
% vim: set spell spelllang=es syntax=tex :
 {\centering \small 
    Introducción a la computación\\
    Departamento de Ingeniería de Computadoras \\
    Facultad de Informática - Universidad Nacional del Comahue \\
    \vspace{20pt} }
\makeatother

\vspace{-2.5cm}
\mbox{\hspace{-1cm}\includegraphics[width=3cm,height=3cm]{logos/uncoma.pdf}\hspace{12cm}
    \includegraphics[width=3cm,height=3cm]{logos/fai.pdf}}



}

% Custom headers and footers
\fancyhf{} % clear all header and footer fields
\fancypagestyle{plain}{\fancyhf{}}
\pagestyle{fancy}
\lhead{\footnotesize Segundo examen parcial}
\rhead{\footnotesize \thepage\ }

\def\ti#1#2{\texttt{#1} & #2 \\ }

\begin{document}

\thispagestyle{empty}
\maketitle
\setlength{\parindent}{1pt}


\begin{enumerate}

    \item Considerando la imagen que se muestra abajo, aplique el esquema de
        codificación de imágenes visto en la práctica. \textbf{No olvide
        incluir la cabecera}. Asuma que 00 = blanco, 11 = negro y 01 = gris.
        Muestre el resultado final expresado en hexadecimal.

    \includegraphics[width=0.50\textwidth]{img/img_2020_1tpo-r.pdf}

    \item Si se tiene una imagen con 200 colores distintos ¿Cuál es la
        cantidad mínima de bits necesarios para la profundidad de color, de
        manera de poder representar la imagen sin pérdida de colores

    \item Dado el siguiente programa:

        \begin{verbatim}
     0: START:  LD  NUM
     1:         JZ  FIN
     2:         ST  OUT
     3:         SUB DOS
     4:         ST  NUM
     5:         JMP START
     6: FIN:    HLT
     7: NUM:    10
     8: DOS:    2
        \end{verbatim}

        \begin{enumerate}

            \item ¿Cuál es la salida del programa?

            \item ¿Cuál es el equivalente en código maquina binario de las
                instrucciones de las direcciones \textbf{0}, \textbf{1} y
                \textbf{5}?

        \end{enumerate}

    \item Dado el siguiente programa en binario:

        \begin{verbatim}
     0: 0100 0010
     1: 0111 1111
     2: 1010 0010
     3: 1111 1101
     4: 0010 0000
     5: 0000 0011
        \end{verbatim}

        ¿Cuál es la salida del programa?

    \item Si se desea extender la funcionalidad de la \emph{MCBE} para que
        tenga un total de 32 instrucciones distintas y la cantidad de celdas
        de memoria se extiende a 256 (incluyendo aquellas mapeadas a la
        \emph{entrada/salida} ¿Cuántos bits deberán tener el código de
        operación y el argumento?

    \item Muestre una secuencia de comandos que permita crear la estructura de
        directorios mostrada en el anexo. Asuma que el directorio raíz (``/'')
        ya existe.

    \item Si el usuario se encuentra posicionado en el directorio
        \emph{notas} del árbol de directorios mostrado en el anexo:

        \begin{enumerate}

            \item Indique una secuencia de comandos para mover el archivo
                \emph{apunte.txt} a la carpeta \emph{linux} utilizando solo
                rutas relativas.

            \item Indique una secuencia de comandos para eliminar el archivo
                \emph{2023.tsv} utilizando solo rutas absolutas.

        \end{enumerate}

\end{enumerate}

\subsection*{ \large\textbf{Anexo:} }

\begin{minipage}{0.45\textwidth}
\textbf{Códigos de operación}

\begin{verbatim}
    010     LD
    011     ST
    100     ADD
    101     SUB
    110     JMP
    111     JZ
    001     HLT
    000     NOP
\end{verbatim}
\end{minipage}
~
\begin{minipage}{0.45\textwidth}
\textbf{Estructura de directorios}

\begin{verbatim}
/
├── notas
│   ├── iadmin
│   │   └── 2024.tsv
│   └── ic
│       ├── 2023.tsv
│       └── 2024.tsv
└── wiki
    ├── c
    │   └── apunte.txt
    └── linux
        ├── ls.2
        └── man.1
\end{verbatim}

\end{minipage}

\end{document}
